\chapter*{Abstract}

This thesis addresses the automation of translating natural-language BDD scenarios into IBDD representations, the first stage of the BDD-to-MBT pipeline. We propose a hybrid approach that combines language-model generation, deterministic syntactic validation through a formal IBDD grammar (Lark/Earley), and an iterative correction loop guided by parser-error diagnostics.

The final evaluation follows a reproducible protocol on the 90-scenario test set, using a fixed dev/test split (30/90), four language configurations (EN--EN, ES--EN, EN--ES, ES--ES), $N=5$ runs per configuration, and \texttt{max-rounds=3}. Results show high final syntactic validity across all configurations (84.22\,\% to 91.11\,\%), with low run-to-run variability. The correction loop improves outcomes in every configuration (gains from 11.78 to 40.22 percentage points) and recovers a substantial fraction of initially invalid translations (53.03\,\% to 73.99\,\%).

The analysis indicates that prompt language is more strongly associated with initial quality than dataset language, while final differences narrow after iterative correction. It also identifies non-convergent cases at the three-round limit, with recurrent failures in expressions and list structures, which defines concrete improvement targets. Overall, the evidence supports the practical syntactic viability of the proposed approach for producing valid IBDD translations under controlled experimental conditions.
